\section{End-to-End Walkthrough}
\label{sec:walkthrough}

This section illustrates the end-to-end behavior of \tool{} using a simplified
transport-parameter example. The purpose is to make the data transformations
concrete before presenting the full implementation details.

\subsection{From Standard Text to a Rule}

Suppose a transport standard states that an endpoint receiving a parameter
below a specified lower bound must treat the message as a protocol error. The
sentence is not enough by itself. The checker must know which field is
constrained, which condition activates the rule, what behavior is required, and
where the requirement appears in the document. \tool{} first resolves the
parameter mention against the field space:
\[
\begin{aligned}
C &: \textit{parameter is present in peer transport parameters},\\
f &: \textit{max\_udp\_payload\_size},\\
A &: \textit{reject values below 1200 with a protocol error},\\
M &: \textit{MUST}.
\end{aligned}
\]
The rule retains the original sentence and section number. If another section
defines the bound or error code, that section is attached as a dependency.

\subsection{From Rule to Code Evidence}

The aligner expands the field name into aliases such as
\texttt{max\_udp\_payload\_size}, \texttt{max\_payload}, and
\texttt{udp\_payload}. Exact matches alone are insufficient: one file may define
the transport-parameter identifier, another may parse the integer, and a third
may validate the configured connection state. \tool{} therefore retrieves a
bundle containing the parser assignment, the validation helper, and the error
return path. Each snippet is tagged with why it was selected: exact field hit,
symbol scope, action-token overlap, or caller expansion.

\subsection{From Evidence to a Judgment}

In the first reasoning round, the model sees the assignment but not the caller
that checks the helper's return value. A one-shot checker may incorrectly report
a missing rejection path. \tool{} instead emits \textsc{Insufficient} with a
typed request for callers of the validation helper. In the second round, the
caller evidence shows that the error return is propagated to connection close.
The final label becomes \textsc{Consistent}. In a different implementation, the
second round might show that the lower-bound check is absent while nearby
upper-bound checks are present, yielding a defect hypothesis.

\subsection{From Hypothesis to Validation}

For a suspected mismatch, the validation agent constructs a parameter block
whose value violates the lower bound. It then chooses the closest available
entry point: a unit-test parser if one exists, otherwise a small harness around
the parser function, otherwise a command-line or integration test. The oracle is
derived from the rule: the implementation should reject the input with the
specified protocol error. The final report records the input, command line,
build options, observed return value, and log excerpt. If the harness cannot
reach the relevant path, the report remains unresolved rather than confirmed.

This walkthrough highlights the main difference between \tool{} and direct LLM
inspection. The model is never asked to solve the entire problem from a large
prompt. It receives a constrained rule, a small evidence bundle, and an explicit
choice to request more evidence when the bundle is incomplete.
